\section{Circuit Complexity --- Karp Lipton Search}
%%% generated by the blueprint pipeline from TCSlib.Complexity.CircuitComplexity.KarpLiptonSearch
\subsection*{Circuit Complexity --- Karp Lipton Search: Introduction}

Declarations of the Lean module \texttt{TCSlib.Complexity.CircuitComplexity.KarpLiptonSearch}.

\subsection*{Circuit Complexity --- Karp Lipton Search: Results}

\begin{definition}[Multi-output directed acyclic circuit]
\lean{BoolCircuit.MultiDAGCircuit}
\label{BoolCircuit.MultiDAGCircuit}
\leanok
\uses{BoolCircuit.DAGGate, BoolCircuit.GatesAcyclic}
For a natural number $n$, a multi-output directed acyclic circuit consists of a list of
gates in topological order, a list of output vertices, a proof that each gate reads only
earlier vertices, and a proof that every listed output is a vertex among the $n$ input
vertices and the gate vertices.
\end{definition}

\begin{definition}[Circuit values on an input]
\lean{BoolCircuit.MultiDAGCircuit.values}
\label{BoolCircuit.MultiDAGCircuit.values}
\leanok
\uses{BoolCircuit.DAGGate.eval, BoolCircuit.runWith}
Given a multi-output directed acyclic circuit with $n$ inputs and an input assignment
$x:\{0,\ldots,n-1\}\to\{\mathrm{false},\mathrm{true}\}$, its values are the list
obtained by running its gates, in order, on the list of input values.
\end{definition}

\begin{definition}[Circuit output evaluation]
\lean{BoolCircuit.MultiDAGCircuit.eval}
\label{BoolCircuit.MultiDAGCircuit.eval}
\leanok
\uses{BoolCircuit.MultiDAGCircuit.values}
Given a multi-output directed acyclic circuit with $n$ inputs and an input assignment
$x$, its evaluation is the list formed, in the specified output order, by reading each
output vertex's value from the circuit's value list; if an index has no value in that
list, the corresponding output is false.
\end{definition}

\begin{definition}[Size of a circuit]
\lean{BoolCircuit.MultiDAGCircuit.size}
\label{BoolCircuit.MultiDAGCircuit.size}
\leanok
For a multi-output directed acyclic Boolean circuit with $n$ inputs and an ordered list
of gates, define its size to be $n$ plus the number of gates.
\end{definition}

\begin{definition}[Fan-in-two circuits]
\lean{BoolCircuit.MultiDAGCircuit.IsFaninTwo}
\label{BoolCircuit.MultiDAGCircuit.IsFaninTwo}
\leanok
A multi-output directed acyclic Boolean circuit is fan-in two if every gate has distinct
input vertices, every NOT gate has exactly one input, and every gate has at most two
inputs.
\end{definition}

\begin{theorem}[Length of circuit evaluation]
\lean{BoolCircuit.MultiDAGCircuit.length_eval}
\label{BoolCircuit.MultiDAGCircuit.length_eval}
\leanok
\uses{BoolCircuit.MultiDAGCircuit.eval}
\difficulty{1}
For every input to a multi-output Boolean circuit, the list of output values has the
same length as the circuit's list of designated output vertices.
\end{theorem}

\begin{definition}[Circuit for one output]
\lean{BoolCircuit.MultiDAGCircuit.outputCircuit}
\label{BoolCircuit.MultiDAGCircuit.outputCircuit}
\leanok
\uses{BoolCircuit.DAGCircuit}
Given a multi-output directed acyclic Boolean circuit with $n$ inputs and an ordered
list of output vertices, an index $i$ smaller than the number of outputs determines a
single-output circuit with the same gates and with the $i$th output vertex as its
output.
\end{definition}

\begin{theorem}[Evaluation of an output circuit]
\lean{BoolCircuit.MultiDAGCircuit.outputCircuit_eval}
\label{BoolCircuit.MultiDAGCircuit.outputCircuit_eval}
\leanok
\uses{BoolCircuit.DAGCircuit.eval, BoolCircuit.DAGCircuit.values, BoolCircuit.MultiDAGCircuit.eval, BoolCircuit.MultiDAGCircuit.outputCircuit, BoolCircuit.MultiDAGCircuit.values}
\difficulty{1}
Let $C$ be a multi-output Boolean circuit, let $i$ be an index less than the number of
its designated output vertices, and let $x$ be an input. The evaluation of the
single-output circuit selecting the $i$th designated output of $C$ equals the $i$th
entry of the evaluation list of $C$.
\end{theorem}

\begin{theorem}[Fan-in of a selected output circuit]
\lean{BoolCircuit.MultiDAGCircuit.outputCircuit_isFaninTwo}
\label{BoolCircuit.MultiDAGCircuit.outputCircuit_isFaninTwo}
\leanok
\uses{BoolCircuit.DAGCircuit.IsFaninTwo, BoolCircuit.MultiDAGCircuit.IsFaninTwo, BoolCircuit.MultiDAGCircuit.outputCircuit}
\difficulty{1}
Let $C$ be a multi-output Boolean DAG circuit with $n$ input vertices, and suppose every
gate reads at most two distinct vertices and every NOT gate reads exactly one vertex.
For any valid output index $i$, the single-output circuit obtained by selecting the
$i$th output has the same property: every gate reads at most two distinct vertices, and
every NOT gate reads exactly one vertex.
\end{theorem}

\begin{definition}[Extension of a relation]
\lean{BoolCircuit.Extends}
\label{BoolCircuit.Extends}
\leanok
Let $R$ be a relation on finite Boolean lists, let $m$ be a natural number, and let $y$
and $p$ be Boolean lists. Say that $p$ extends to an $R$-related list of length $m$
relative to $y$ if there is a Boolean list $s$ such that $p$ followed by $s$ has length
$m$ and is related to $y$ by $R$.
\end{definition}

\begin{definition}[Recursive search bits]
\lean{BoolCircuit.searchBits}
\label{BoolCircuit.searchBits}
\leanok
\uses{BoolCircuit.DAGCircuit.eval}
Let $n$ be a natural number and let $D_j$ be a Boolean circuit with $n+j$ inputs for
each natural number $j$. For a Boolean list $y$, define a list of search bits
recursively: the list at stage $0$ is empty, and the list at stage $j+1$ consists of the
stage-$j$ list followed by the output of $D_{j+1}$ on the inputs given by $y$, the
stage-$j$ search bits, and then a final true bit.
\end{definition}

\begin{theorem}[Length of the search bits]
\lean{BoolCircuit.length_searchBits}
\label{BoolCircuit.length_searchBits}
\leanok
\uses{BoolCircuit.searchBits}
\difficulty{4}
Let $D$ be a family of directed acyclic Boolean circuits indexed by natural numbers,
where the circuit at index $j$ has $n+j$ inputs. For every Boolean list $y$ and natural
number $j$, the list of search bits constructed from $y$ using the first $j$ circuits
has length $j$.
\end{theorem}

\begin{theorem}[Correctness of sequential bit search]
\lean{BoolCircuit.searchBits_spec}
\label{BoolCircuit.searchBits_spec}
\leanok
\uses{BoolCircuit.DAGCircuit.eval, BoolCircuit.Extends, BoolCircuit.length_searchBits, BoolCircuit.searchBits}
\difficulty{6}
Let $n,m$ be natural numbers, let $R$ be a relation between Boolean lists, and let
$(D_j)_{j\in\mathbb N}$ be a family of Boolean circuits, with $D_j$ taking $n+j$ input
bits. Suppose that for every $1\le j\le m$, every Boolean list $y$ of length $n$, and
every Boolean list $p$ of length $j$, the circuit $D_j$ outputs true on the
concatenation of $y$ and $p$ exactly when some extension of $p$ to a list of length $m$
is related to $y$ by $R$. If $y$ has length $n$ and some list of length $m$ is related
to $y$ by $R$, then the list produced by the sequential search using these circuits is
related to $y$ by $R$.
\end{theorem}

\begin{definition}[Search state update]
\lean{BoolCircuit.searchStep}
\label{BoolCircuit.searchStep}
\leanok
\uses{BoolCircuit.DAGGate, BoolCircuit.embedGates}
Let $n$ be a natural number and let $D_j$ be a directed acyclic Boolean circuit with
$n+j$ inputs for each natural number $j$. Given a natural number $j$ and a state
consisting of a list of gates and a list of vertex indices, append to the gate list the
embedded gates of $D_{j+1}$, using as sources the indices $0,\ldots,n-1$, the existing
vertex indices, and then $n$, with the embedding offset $n$ plus the current gate-list
length. Append to the index list the vertex index $n$ plus the current gate-list length
plus the number of gates of $D_{j+1}$.
\end{definition}

\begin{definition}[Recursive Boolean circuit search state]
\lean{BoolCircuit.searchState}
\label{BoolCircuit.searchState}
\leanok
\uses{BoolCircuit.DAGGate, BoolCircuit.constGate, BoolCircuit.searchStep}
Let $n$ be a natural number, and for each natural number $j$ let $D_j$ be a directed
acyclic Boolean circuit with $n+j$ input vertices. Define $S_0$ to be the pair
consisting of the single true constant gate and the empty list of indices. Recursively,
if $S_j=(G,I)$, form a circuit gate list from $D_{j+1}$ by mapping each argument $v$ to
the $(v+1)$st entry of the list formed by concatenating $0,1,\ldots,n-1$, the entries of
$I$, and $n$ when $v<n$, and to $n+|G|+(v-n)$ otherwise; remove repeated arguments from
each mapped gate, then append an AND gate whose sole argument is the mapped output of
$D_{j+1}$. The state $S_{j+1}$ is the pair whose first list is $G$ followed by this
constructed gate list, and whose second list is $I$ followed by the index
$n+|G|+|\mathrm{gates}(D_{j+1})|$.
\end{definition}

\begin{theorem}[Search state gate bound]
\lean{BoolCircuit.length_searchState_le}
\label{BoolCircuit.length_searchState_le}
\leanok
\uses{BoolCircuit.DAGCircuit.size, BoolCircuit.length_embedGates, BoolCircuit.searchState, BoolCircuit.searchStep}
\difficulty{4}
Let $D$ be a family of directed acyclic circuits, with the circuit at index $j$ having
$n+j$ inputs, and let $S,M$ be natural numbers. Suppose every index $j$ with $1\le j\le
M$ has circuit size at most $S$. Then, for every $j\le M$, the gate list in the state
after $j$ search steps has length at most $1+j(S+1)$.
\end{theorem}

\begin{theorem}[Structure of the search state]
\lean{BoolCircuit.searchState_struct}
\label{BoolCircuit.searchState_struct}
\leanok
\uses{BoolCircuit.DAGCircuit.IsFaninTwo, BoolCircuit.GatesAcyclic, BoolCircuit.GatesAcyclic.append, BoolCircuit.GatesAcyclic.nil, BoolCircuit.GatesAcyclic.snoc, BoolCircuit.constGate, BoolCircuit.faninTwo_embedGates, BoolCircuit.gatesAcyclic_embedGates, BoolCircuit.length_embedGates, BoolCircuit.searchState, BoolCircuit.searchStep}
\difficulty{5}
Let $n,j\in\mathbb{N}$, and for each $r\in\mathbb{N}$ let $D_r$ be a Boolean circuit on
$n+r$ inputs in which every gate has distinct arguments, every negation gate has exactly
one argument, and every gate has at most two arguments. After $j$ stages of the search
construction, the gate list is acyclic, its output list has length $j$ and contains only
vertices of the circuit, and the vertex $n$ is a vertex of the circuit. Every gate in
the list has distinct arguments, every negation gate has exactly one argument, and every
gate has at most two arguments.
\end{theorem}

\begin{theorem}[Values produced by the search construction]
\lean{BoolCircuit.searchState_values}
\label{BoolCircuit.searchState_values}
\leanok
\uses{BoolCircuit.DAGCircuit.IsFaninTwo, BoolCircuit.DAGCircuit.eval, BoolCircuit.DAGGate.eval, BoolCircuit.embedGates, BoolCircuit.getD_map_getD, BoolCircuit.runWith, BoolCircuit.runWith_append, BoolCircuit.runWith_embedGates_getD, BoolCircuit.runWith_getD_of_lt, BoolCircuit.runWith_singleton, BoolCircuit.searchBits, BoolCircuit.searchState, BoolCircuit.searchState_struct, BoolCircuit.searchStep}
\difficulty{6}
Let $n,j\in\mathbb{N}$, let $x$ be a Boolean assignment to the $n$ input bits, and for
each $r\in\mathbb{N}$ let $D_r$ be a Boolean circuit on $n+r$ inputs in which every gate
has distinct arguments, every negation gate has exactly one argument, and every gate has
at most two arguments. Define $b_0$ to be the empty list and, for $r\ge1$, define $b_r$
by appending to $b_{r-1}$ the value of $D_r$ on the input consisting of $x$, the bits of
$b_{r-1}$, and a final true bit. Evaluating the gates of the search construction after
$j$ stages on $x$, the value at vertex $n$ is true, and the values at its listed output
vertices, in order, are exactly $b_j$.
\end{theorem}

\begin{definition}[Circuit built from the search state]
\lean{BoolCircuit.searchCircuit}
\label{BoolCircuit.searchCircuit}
\leanok
\uses{BoolCircuit.DAGCircuit.IsFaninTwo, BoolCircuit.MultiDAGCircuit, BoolCircuit.searchState, BoolCircuit.searchState_struct}
Let $n$ be a natural number, and for each natural number $j$ let $D_j$ be a directed
acyclic Boolean circuit with $n+j$ inputs. Suppose every $D_j$ has distinct input
vertices at each gate, every NOT gate has exactly one input, and every gate has at most
two inputs. For each natural number $m$, form a multi-output directed acyclic Boolean
circuit with $n$ inputs whose gate list and ordered output list are the two components
of the search state at stage $m$.
\end{definition}

\begin{theorem}[Evaluation of the search circuit]
\lean{BoolCircuit.searchCircuit_eval}
\label{BoolCircuit.searchCircuit_eval}
\leanok
\uses{BoolCircuit.DAGCircuit.IsFaninTwo, BoolCircuit.MultiDAGCircuit.eval, BoolCircuit.searchBits, BoolCircuit.searchCircuit, BoolCircuit.searchState_values}
\difficulty{1}
Let $n,m$ be natural numbers, and let $(D_j)_{j\in\mathbb N}$ be a family of Boolean
circuits, with $D_j$ taking $n+j$ input bits. If every $D_j$ is well formed and has
fan-in at most two, then for every input $x\in\{0,1\}^n$, the search circuit built from
the family outputs exactly the list produced by the sequential bit search on the bit
list of $x$.
\end{theorem}

\begin{theorem}[Fan-in bound for the search circuit]
\lean{BoolCircuit.searchCircuit_isFaninTwo}
\label{BoolCircuit.searchCircuit_isFaninTwo}
\leanok
\uses{BoolCircuit.DAGCircuit.IsFaninTwo, BoolCircuit.MultiDAGCircuit.IsFaninTwo, BoolCircuit.searchCircuit, BoolCircuit.searchState_struct}
\difficulty{3}
Let $n,m\in\mathbb{N}$, and for each $j\in\mathbb{N}$ let $D_j$ be a Boolean circuit on
$n+j$ inputs in which every gate has distinct arguments, every negation gate has exactly
one argument, and every gate has at most two arguments. Then the circuit with multiple
outputs produced by the search construction from $D_0,D_1,\ldots$ after $m$ stages has
distinct arguments at every gate, exactly one argument at each negation gate, and at
most two arguments at every gate.
\end{theorem}

\begin{theorem}[Size bound for the search circuit]
\lean{BoolCircuit.searchCircuit_size_le}
\label{BoolCircuit.searchCircuit_size_le}
\leanok
\uses{BoolCircuit.DAGCircuit.IsFaninTwo, BoolCircuit.DAGCircuit.size, BoolCircuit.MultiDAGCircuit.size, BoolCircuit.length_searchState_le, BoolCircuit.searchCircuit}
\difficulty{3}
Let $n,m,S\in\mathbb{N}$, and for each $j\in\mathbb{N}$ let $D_j$ be a Boolean circuit
on $n+j$ inputs in which every gate has distinct arguments, every negation gate has
exactly one argument, and every gate has at most two arguments. Suppose that for every
$j$ with $1\le j\le m$, the size of $D_j$ is at most $S$. Then the size of the
multiple-output circuit produced by the search construction from this family after $m$
stages is at most $n+1+m(S+1)$.
\end{theorem}

\begin{theorem}[Correctness of the search circuit]
\lean{BoolCircuit.searchCircuit_correct}
\label{BoolCircuit.searchCircuit_correct}
\leanok
\uses{BoolCircuit.DAGCircuit.IsFaninTwo, BoolCircuit.DAGCircuit.eval, BoolCircuit.Extends, BoolCircuit.MultiDAGCircuit.eval, BoolCircuit.length_searchBits, BoolCircuit.searchBits_spec, BoolCircuit.searchCircuit, BoolCircuit.searchCircuit_eval}
\difficulty{3}
Let $n,m$ be natural numbers, let $R$ be a relation between Boolean lists, and let
$(D_j)_{j\in\mathbb N}$ be a family of Boolean circuits, with $D_j$ taking $n+j$ input
bits. Suppose every $D_j$ is well formed, has fan-in at most two, and, for each $1\le
j\le m$, every Boolean list $y$ of length $n$ and list $p$ of length $j$ satisfy: $D_j$
outputs true on their concatenation exactly when $p$ has an extension of length $m$
related to $y$ by $R$. For any input $x\in\{0,1\}^n$, if there is a length-$m$ list
related to the list of bits of $x$ by $R$, then the search circuit built from this
family outputs a list of length $m$ that is related to the list of bits of $x$ by $R$.
\end{theorem}
